\icode{E-MK-SI-L3-260625-001}
图像法、微元累积法是常用的物理研究方法。如图所示某质点沿直线运动，运动速率的倒数\(\frac{1}{v}\)与位移\(x\)的关系（\(OA\)与\(AA^\prime\)距离相等），关于质点的运动，下列说法正确的是\paren[]

\begin{tikzpicture}[>=stealth,font=\scriptsize]
	\draw[->,blue] (0,0)--(0,3)node[right]{\(\frac{1}{v}\)/(\qty{}{s\cdot m^{-1}})};
    \draw[->,blue] (0,0)--(3,0)node[right]{\(x\)/\qty{}{m}};
    \draw[purple] (0,0)--(2.2,2.2);
    \draw[dashed,cyan] (0,.8)node[left]{\(A\)}--(.8,.8)--(.8,0)node[below]{\(C\)};
    \draw[dashed,cyan] (0,1.6)node[left]{\(A^\prime\)}--(1.6,1.6)--(1.6,0)node[below]{\(C^\prime\)};
    \node at(-.2,-.2) {\(O\)};
\end{tikzpicture}
\begin{choices}
    \item 质点做匀变速直线运动
    \item \(\frac{1}{v} - x\)图线斜率等于质点运动的加速度
    \item \(\frac{1}{v} - x\)图线与轴围成的面积没有实际的物理意义
    \item 质点从\(C\)运动到\(C^\prime\)所用的运动时间是从\(O\)运动到\(C\)所用时间的3倍
\end{choices}
 
